\section{Deninger's comparison map and the trace formulas in the letter}
\label{sec:ACtracefollowup}

Here \(\mathrm{M25}\) denotes M.~Morishita,
\emph{On a relation between Deninger's foliated dynamical systems and
Connes--Consani's adelic spaces}, arXiv:2508.15971, supplied file
\texttt{Deninger\_Connes-Consani.tex}.  Its equation numbers are printed
explicitly in the source, rather than generated by LaTeX labels.
\(\mathrm{L26}\) denotes A.~Connes,
\emph{The Riemann Hypothesis: Past, Present and a Letter Through Time},
arXiv:2602.04022, supplied file \texttt{rhready.tex}.
We retain the definitions and normalizations of AC-T01--AC-T14.
The theorem \(A\) remains a stipulated input.

\subsection{A necessary correction to the character space}

\begin{proposition}[AC-T15: the finite-kernel condition cannot be omitted]
\label{AC-T15}
Fix a rational prime \(p\) and an isomorphism
\(\iota:\overline{\mathbb F}_p^\times\longrightarrow\mu_{(p)}
\subset\mathbb C^\times\).  There is an isomorphism
\[
 \operatorname{Hom}(\overline{\mathbb F}_p^\times,\mathbb C^\times)
   \simeq \prod_{\ell\ne p}\mathbb Z_\ell,\qquad
 b\longmapsto\bigl(z\longmapsto\iota(z)^b\bigr).
\]
Writing \(U_{(p)}=\prod_{\ell\ne p}\mathbb Z_\ell^\times\), the image of
\[
 U_{(p)}\times\mathbb N\longrightarrow
 \operatorname{Hom}(\overline{\mathbb F}_p^\times,\mathbb C^\times),
 \qquad (u,n)\longmapsto\iota^{nu}
\]
is exactly the set of characters with finite kernel.  In particular,
this map is not surjective onto the unrestricted character group.

Thus M25, (2.2.7), requires a finite-kernel restriction.
Remark 2.1.13 explicitly omits the admissibility conditions, while
(2.1.5) uses unrestricted homomorphisms.  Under those literal
definitions, the asserted surjection and the ensuing description of
the entire packet in Theorems 2.2.8--2.2.9 cannot be used as stated.
\end{proposition}
\begin{proof}
The group \(\overline{\mathbb F}_p^\times\) is the direct sum of its
Pruefer \(\ell\)-groups for \(\ell\ne p\).  A homomorphism from its
\(\ell\)-primary part to \(\mathbb C^\times\) has image in the same
\(\ell\)-primary roots of unity.  Compatible exponentiation maps on
the cyclic groups \(\mu_{\ell^r}\) are precisely elements of
\(\varprojlim_r\mathbb Z/\ell^r\mathbb Z=\mathbb Z_\ell\).
Taking the product over \(\ell\ne p\) proves the first assertion.

If \(b_\ell\ne0\), multiplication by \(b_\ell\) on the Pruefer
\(\ell\)-group has kernel of order \(\ell^{v_\ell(b_\ell)}\).
If \(b_\ell=0\), that kernel is infinite.  Consequently the full
kernel is finite exactly when every \(b_\ell\) is nonzero and
\(v_\ell(b_\ell)=0\) for all but finitely many \(\ell\).
In that event set
\[
 n=\prod_{\ell\ne p}\ell^{v_\ell(b_\ell)}.
\]
Then \(b=nu\) with \(u\in U_{(p)}\); conversely every such product
has finite kernel.  The trivial character corresponds to \(b=0\)
and is already a counterexample to the claimed unrestricted
surjectivity.  Another counterexample is \(b_\ell=\ell\) for every
\(\ell\ne p\), whose kernel contains a nontrivial subgroup of every
prime order other than \(p\).
\end{proof}

This correction agrees with the actual restrictions in
C.~Deninger, \emph{Dynamical systems for arithmetic schemes},
arXiv:1807.06400v4, \S7, equations (62)--(66): the subsets
\(H_{\rm fin,p}\) require finite kernel away from \(p\).
Equation (63) identifies them with
\(\mathbb N(U_{(p)}\times0)\).
Only that stated restriction and those displayed identifications
are used here; no claim is made to have checked every proof in
Deninger's 119-page paper.

\subsection{The continuous map survives, with a corrected orbit scope}

\begin{proposition}[AC-T16: the map and its admissible prime orbits]
\label{AC-T16}
Let \(F/\mathbb Q\) be finite abelian and write
\(U_F=\operatorname{Gal}(\mathbb Q^{\rm ab}/F)\).
On the rational-Witt character spaces of M25, the map
\[
 \Psi_F:\mathfrak X_F\longrightarrow
 \mathscr X_F=\mathbb Q^\times\backslash\mathbb A_{\mathbb Q}/U_F
\]
constructed in Lemmas 3.4--3.5 and Theorem 3.6(1) is continuous,
Galois equivariant, and reverses the multiplicative flow:
\[
 \Psi_F(\phi^r x)=\phi^{r^{-1}}\Psi_F(x).
\]
These assertions hold before imposing the finite-kernel condition.
On the finite-kernel part over a closed point of residue
characteristic \(p\), its finite adelic coordinate has zero
\(p\)-component and is a finite idele at the other primes.
For \(F=\mathbb Q\), each of the resulting periodic orbits maps
onto the adelic prime circle \(C_p\).
For finite abelian \(F\) in which \(p\) is unramified, the finite
cover over such an orbit has the usual Artin monodromy and splits
into \([F:\mathbb Q]/f_p\) circles of length \(f_p\log p\).
All of these statements are unconditional.
\end{proposition}
\begin{proof}
At the \(\mathbb Q^{\rm ab}\) level, restrict a multiplicative
character \(P\) to the roots of unity after reduction at its prime.
This gives an element \(\psi(P)\in
\operatorname{End}(\mu_\infty)=\widehat{\mathbb Z}\).
Restriction is continuous for pointwise convergence.
The relations \(F_n[\zeta]=[\zeta^n]\) and
\(\sigma(\zeta)=\zeta^{\operatorname{cycl}(\sigma)}\) give
\[
 \psi(F_nP)=n\psi(P),\qquad
 \psi(P^\sigma)=\operatorname{cycl}(\sigma)\psi(P).
\]
Passing to the localization
\(\widehat{\mathbb Z}\otimes\mathbb Q=\mathbb A_{\mathbb Q,f}\)
therefore gives a continuous map \(\check\psi\) with the
same equivariance for positive rational scalars.  On the suspended
space define
\[
 [(P,r)]\longmapsto
 \mathbb Q^\times\bigl(\check\psi(P),r^{-1}\bigr).
\]
Replacing \((P,r)\) by \((F_qP,q^{-1}r)\) multiplies both coordinates
of its image by \(q\), so the map is well defined.  The displayed
Galois relation descends through \(U_F\), and replacing \(r\) by
\(tr\) replaces the real coordinate by \(t^{-1}r^{-1}\).
This proves the continuous equivariance assertions directly.

At a prime over \(p\), reduction kills the \(p\)-power roots of
unity, so the \(p\)-component is zero.  On the finite-kernel
part, AC-T15 shows that the other components are \(nu\), with
\(n\in\mathbb N\) and \(u\in U_{(p)}\), before localization.
After multiplying by a rational scalar they are units.  Over
\(\mathbb Q\), quotienting by \(\widehat{\mathbb Z}^{\times}\)
leaves the representative with all finite components equal to
one except for the zero at \(p\).  Its positive-real orbit is
the circle with stabilizer \(p^{\mathbb Z}\), and the flow map
is onto.

For the finite unramified covering, put
\(G=\operatorname{Gal}(F/\mathbb Q)\) and
\(g_p=\operatorname{Frob}_p\in G\).
The stabilizer of a finite-kernel residue character in the
abelian decomposition group is inertia: a Frobenius element
fixing it must act trivially after multiplication by each
nonzero \(\ell\)-adic component in AC-T15, hence acts trivially
on every prime-to-\(p\) root of unity.  Those roots generate
\(\overline{\mathbb F}_p\), so its residue-field action is the
identity.  For the unramified extension \(F\), inertia maps
trivially to \(G\).  Thus the fiber of the finite cover is a
\(G\)-torsor, and a trip around the prime circle acts by
\(g_p\).  This proves, directly on the restricted character
space, the suspension description used next.
The covering is the suspension of the permutation
\(g\mapsto g_pg\) of \(G\); equivalently it is the quotient of
\(G\times\mathbb R\) by
\((g,t)\sim(g_pg,t-\log p)\).
Each cycle of this permutation has length
\(f_p=\operatorname{ord}(g_p)\).  Gluing its \(f_p\) intervals
gives a circle of length \(f_p\log p\); there are
\(|G|/f_p\) such cycles.  This also proves the asserted
monodromy description.  Reversing flow orientation replaces
\(g_p\) by \(g_p^{-1}\); one must keep this reversal when
comparing the two suspensions via \(\Psi_F\).
\end{proof}

The \(p\)-component calculation alone, used in M25's proof of
Theorem 3.6(2), does not prove that the other components form a
finite idele.  AC-T15 supplies exactly the omitted condition.
The unrestricted space has points whose image has additional
zero components, as well as points with infinitely many nonunit
components.  No assertion that all such points lie on a prime
circle is made here.

\begin{proposition}[AC-T17: the infinite fiber is a solenoid]
\label{AC-T17}
Put \(H_p=\overline{p^{\mathbb Z}}\subset U_{(p)}\).
For the mapping torus
\[
 M_p=U_{(p)}\times_{p^{\mathbb Z}}\mathbb R_{>0},
\]
the fiber over \(aH_p\) of the natural map
\(M_p\to U_{(p)}/H_p\) is
\[
 H_p\times_{p^{\mathbb Z}}\mathbb R_{>0}.
\]
It is compact and its positive-real flow orbits are injective,
dense, and nonclosed.  In particular the fiber is not
homeomorphic to \(\mathbb R_{>0}\).
This corrects the fiber assertion in M25, Theorem 2.2.8(2),
and the corresponding infinite-cover assertion in Theorem 2.3.1,
even after finite kernels have been restored.
\end{proposition}
\begin{proof}
The inverse image of \(aH_p\) in \(U_{(p)}\) is \(aH_p\);
translation by \(a^{-1}\) identifies its suspension with the
displayed fiber.  A fundamental domain is
\(H_p\times[1,p]\), with its two compact ends glued by multiplication
by \(p\).  Thus the fiber is compact.
In logarithmic coordinates its flow orbit through \([h,0]\)
returns to the transversal at the points \(hp^n\).
These are dense in \(hH_p=H_p\) by definition, so the flow orbit
is dense in the fiber.
A period would require \(p^n=1\) in \(U_{(p)}\) for a nonzero
integer \(n\).  The diagonal embedding into any
\(\mathbb Q_\ell\), \(\ell\ne p\), is injective on
\(\mathbb Q^\times\), so this is impossible.
The orbit cannot be all of the fiber: its intersections with
the transversal are countable, whereas \(H_p\) is an infinite
profinite group and is uncountable.  Indeed a countable compact
Hausdorff group would have an isolated point by the Baire
theorem, and translation would then make it discrete and
finite.  Hence the dense orbit is
not closed.  Compactness also directly excludes the asserted
homeomorphism with \(\mathbb R_{>0}\).
\end{proof}

\subsection{What \(A\) adds to the abelian orbit data}

\begin{theorem}[AC-T18: abelian orbit factors and the exact analytic consequence]
\label{AC-T18}
Let \(F/\mathbb Q\) be finite abelian and \(G\) its Galois group.
At every rational prime \(p\), write \(I_p\) for inertia and let
\(\operatorname{Frob}_p\) act on the inertia invariants of the
regular representation.  Then
\[
 \det\bigl(1-p^{-s}\operatorname{Frob}_p\mid
                    \mathbb C[G]^{I_p}\bigr)^{-1}
   =\prod_{\mathfrak p\mid p}(1-(N\mathfrak p)^{-s})^{-1}.
\]
Decomposing by characters gives the primitive Dirichlet
factorization
\[
 \zeta_F(s)=\prod_{\chi\in\widehat G}L(s,\chi).
\]
Under \(A\), all its nontrivial zeros have real part \(1/2\),
with multiplicities given by the sum of the multiplicities of
the factors.  Every character's centered Weil form is positive,
and the finite-window positive operators and diagonal spectral
realizations of AC-T04--AC-T06 and AC-T13 apply to it.

The comparison map \(\Psi_F\) is a map of topological dynamical
spaces.  The foregoing argument constructs no intertwiner
between leafwise cohomology and the Hilbert trace spaces, and
proves no new Lefschetz trace formula on the unrestricted
Deninger space.
\end{theorem}
\begin{proof}
The invariant subspace has the permutation basis indexed by
\(G/I_p\).  The Frobenius permutation has one cycle for each
prime of \(F\) above \(p\), of length its residue degree \(f_p\).
The determinant of \(1-zC_f\) for a cyclic permutation of length
\(f\) is \(1-z^f\), proving the local formula, including ramified
primes.  The regular representation of the abelian group is the
direct sum of its one-dimensional characters.  Its \(I_p\)
invariants retain precisely those trivial on inertia.
Their Euler factors are precisely the primitive Dirichlet
factors; ramified characters contribute one.
Multiplication of the absolutely convergent Euler products on
\(\operatorname{Re}s>1\), followed by their known meromorphic
continuations, proves the factorization.

Inside \(0<\operatorname{Re}s<1\) the factors are holomorphic,
so zeros of the product are exactly their zero multisets with
summed orders.  The input \(A\) puts each such multiset on the
critical line.  Applying the proved formulas in AC-T02,
AC-T04--AC-T06 and AC-T13 yields the remaining conclusions.
The construction just used only a finite permutation
representation, its character decomposition, and analytic
continuation of classical \(L\)-functions.  It neither defines
nor uses a map on a cohomology group of \(\mathfrak X_F\).
\end{proof}

The precise remaining geometric bridge would need a specified
cohomology theory or Hilbert space, a domain for its flow generator,
a trace or regularized trace, and an intertwining result preserving
local multiplicities and the archimedean terms.  The reversal of
flow in AC-T16 also reverses the centered spectral parameter.
This lists missing hypotheses in the proposed construction; it is
not a claim that their truth is logically independent of \(A\).
The finite-dimensional character constituents above are ordinary
Artin representations.  Their finite-place Frobenius eigenvalues
are roots of unity.  The location of the global analytic zeros,
and a positive Hilbert realization of those zeros, are different
statements from that elementary local weight-zero property.

\subsection{The two-cutoff trace formula in the letter}

\begin{proposition}[AC-T19: the archimedean identity is an ordinary trace]
\label{AC-T19}
On \(L^2(\mathbb R)_{\rm ev}\), let \(P_T\) be multiplication by
\(1_{[-T,T]}\), \(Q_W=\mathcal F P_W\mathcal F^{-1}\), and let
\(\vartheta\) be the unitary dilation representation used in L26.
For every \(T,W>0\) and \(f\in C_c^\infty(\mathbb R_{>0})\),
the operator
\(\vartheta(f)(1-P_T-Q_W)\) is trace class and
\[
 W_\infty(f)=\log(TW)f(1)
       +\operatorname{Tr}\bigl(\vartheta(f)(1-P_T-Q_W)\bigr).
 \tag{TF.1}
\]
This is L26, equation \texttt{weilq}, and is unconditional.
It is not the global prolate-image cutoff limit of C99,
Theorem VIII.5 and the final paragraphs of \S VIII.
\end{proposition}
\begin{proof}
We first set \(T=W=1\).
Use the unitary logarithmic coordinate on even functions,
and let \(I\) be reflection in that coordinate.
If \(P_{\rm out}=1-P_1\) and \(V=I\mathcal F\), then
\[
 V^*P_{\rm out}V=Q_1,\qquad
 -V^*[P_{\rm out},V]=1-P_1-Q_1.
 \tag{TF.2}
\]
The first identity uses \(IP_{\rm out}I=P_1\) and the fact that
the Fourier transform is an involution on the even subspace.
In Mellin coordinates \(V\) is multiplication by the
archimedean gamma phase.  CC20, Proposition \texttt{lemsch7}(iv),
equation \texttt{traceequa}, identifies
\(\operatorname{Tr}(\vartheta(f)V^*[P_{\rm out},V])\)
with \(W_{\mathbb R}(f)=-W_\infty(f)\) and proves trace class.
For clarity, its trace-class argument is that the gamma phase
and all its derivatives have at most polynomial growth, while
the Mellin transform of \(f\) is Schwartz.  The product rule
writes its smoothed commutator as a difference of two
commutators with Schwartz multiplication symbols.  CC20,
Appendix A, Lemma \texttt{quantsmooth}, proves these trace
class by the rapidly decreasing Hankel kernel.  Together
with (TF.2), this proves (TF.1) at \(T=W=1\).

In logarithmic coordinates \(P_T-P_1\) is multiplication by
a signed bounded interval of oriented length \(\log T\).
Convolution by the smooth compactly supported logarithmic
function \(f\), followed by multiplication by a bounded
interval, is trace class.  Indeed, applying any power of the
output harmonic oscillator to its kernel still gives a
square-integrable kernel; factoring through a sufficiently
large inverse power of that oscillator gives a product of
two Hilbert--Schmidt operators.  Compressing to the interval
and integrating the continuous diagonal kernel gives
\[
 \operatorname{Tr}\bigl(\vartheta(f)(P_T-P_1)\bigr)
     =f(1)\log T.
 \tag{TF.3}
\]
The same formula for \(Q_W-Q_1\) follows by Fourier conjugation:
Fourier transforms dilation by \(r\) into dilation by \(r^{-1}\),
and the inverted test has the same value at \(1\).
Subtracting these two trace-class differences from (TF.2)
proves (TF.1) for arbitrary \(T,W\).
\end{proof}

\begin{proposition}[AC-T20: the semilocal positive trace and its exact limitation]
\label{AC-T20}
Take the \(K_S\)-invariant semilocal Hilbert model in L26,
equations \texttt{wS}--\texttt{Modulus}, and its two module cutoffs.
The formula stated there as \texttt{weilqb} is
\[
 -\sum_{v\in S}W_v(f)
   =\log(TW)f(1)+
      \operatorname{Tr}\bigl(\vartheta(f)B^S_{T,W}\bigr),
 \qquad B^S_{T,W}=1-P^S_T-\widehat P^S_W.
 \tag{TF.4}
\]
It is the semilocal trace identity, independent of \(A\).
For \(v\in C_c^\infty(-L,L)\), \(g(e^t)=v(t)\),
\(f=g*g^*\), and \(S\) containing every prime \(p\le e^{2L}\),
the input \(A\) gives the precise consequence
\[
 \operatorname{Tr}\bigl(\vartheta(f)B^S_{T,W}\bigr)
 +\log(TW)\|v\|_2^2
 +2\operatorname{Re}\bigl(\ell_+(v)\overline{\ell_-(v)}\bigr)
 =q_{1,e^L}(v)\ \geq c_L\|v\|_2^2,
 \quad c_L>0.
 \tag{TF.5}
\]
Here \(c_L\) is the lowest eigenvalue in AC-T13, and
\(\ell_\pm(v)=\int v(t)e^{\pm t/2}\,dt\).
In particular \(TW=1\) and either vanishing moment give a
strictly positive trace for every nonzero such \(v\).

Positivity of these traces is not positivity of the bare
two-cutoff operator.  Already in the archimedean case,
\(B_{T,W}=1-P_T-Q_W\) has a negative eigenvalue for every
\(T,W>0\).
\end{proposition}
\begin{proof}
The identification of (TF.4) with the semilocal formula is the
one expressly stated after L26's equation \texttt{Modulus};
its \(K_S\)-invariant model is essential.  The new consequence
does not replace that model by the global image projection
of C99.  Since \(f\) is supported in
\([e^{-2L},e^{2L}]\), every nonzero prime-power term belongs
to \(S\).  Also \(f(1)=\|v\|_2^2\).  Thus AC-T05, with its
two pole terms retained, turns (TF.4) into the equality in
(TF.5); AC-T13 proves its strict lower bound.

For the last assertion, the operator \(P_TQ_WP_T\) on the
time interval is nonzero, positive, and compact: its
frequency kernel is continuous on the compact square.
It has an eigenvector \(e\), normalized, with eigenvalue
\(\sigma^2>0\).  Necessarily \(\sigma<1\), since a common
time- and frequency-supported nonzero function would have
an entire Fourier transform vanishing on an open set,
which is impossible.  Put \(h=Q_We/\sigma\).
Then \(\|h\|=1\), \(P_Th=\sigma e\), and \(Q_We=\sigma h\).
Consequently
\[
 (1-P_T-Q_W)(e+h)=-\sigma(e+h).
 \tag{TF.6}
\]
This is a nonzero negative eigenvector.  The calculation is
unconditional and involves the actual cutoff projections.
\end{proof}

The trace assertion (TF.4) is used as the semilocal identity
stated in L26 and C99.  AC-T19 supplies a full verification
of its archimedean two-parameter specialization.  No
new independent proof of all of C99's semilocal theorem is
claimed here.  The separate global gate T-CUT in AC-T11
concerns prolate \emph{image} projections and their
conductor-dependent Gram matrices; (TF.1)--(TF.5) do not
identify those projections.

\subsection{Status inventory for these added sources}

M25's Lemmas 3.4--3.5 and Theorem 3.6(1) give the continuous
comparison, proved again in AC-T16.
Its Theorems 1.2.4 and 1.2.7 describe adelic mapping tori and
finite Artin monodromy; the finite permutation computation
in AC-T16 and AC-T18 is sufficient for the uses made here.
Propositions 2.1.1 and 2.1.9 supply rational-Witt descent and
are background theorems cited there from Deninger and
Kucharczyk--Scholze.  They are not new consequences of \(A\).
The unrestricted assertions arising from (2.2.7) require the
correction AC-T15, and the infinite-fiber assertions require
AC-T17.  The introduction describes Deninger's
cohomological/Lefschetz programme as conjectural; no such
cohomology or trace comparison is established by Theorem 3.6.

L26's equations \texttt{weilq} and \texttt{weilqb} are
unconditional trace identities, with the new positive
specialization (TF.5) supplied by \(A\).
Its sheaf theorem \texttt{structure1intro} is an existing
construction, not a result requiring RH.
Its conditionally real Fourier zeros in \texttt{mainintro},
the approximate-kernel limit \texttt{hermfact1}, the two
remaining steps in \texttt{miss}, and the RH-conditional heat
expansion \texttt{heatkernelintro} are treated in the spectral
module and its followup.  The semilocal identity in
\texttt{sectsemiloc} does not remove the eigenvector comparison
explicitly retained in \texttt{miss}.
